% ============================================================================
% 02_architecture | 06_attention
% Section-body file: include immediately after the relevant \section command
% in the master tstextbook.tex file. This file intentionally contains no
% document class, package imports, \chapter, \section, or bibliography command.
% Figure sources are TeX fragments stored under the root-level figures folder.
% All citations refer only to the papers supplied for this section.
% ============================================================================

The preceding discussion examined normalization as a means of controlling the
scale and statistical behavior of intermediate representations. A further
architectural question remains: how can the representation of one token
incorporate information carried by other tokens in the same sequence? Token
embeddings and positional information establish the initial numerical
representations, while normalization and position-wise transformations operate
on those representations. None of these operations, by itself, specifies a
content-dependent mechanism through which one position can retrieve
information from another. The attention mechanism supplies this missing
operation. It computes relationships between token representations and uses
those relationships to determine how information is exchanged across sequence
positions. Attention is therefore the next step in the architectural account
of how a language model converts a sequence of token representations into
context-dependent representations
\parencite{ansar2024transformers,dao2023flashattention2}.

\subsection{The role of attention in sequence representation}
\label{subsec:attn-role-sequence-representation}

Consider a sequence containing \(N\) tokens. After tokenization, embedding,
and the addition or incorporation of positional information, each token is
represented by a vector in a \(d_{\mathrm{model}}\)-dimensional real vector
space. Using a row-vector convention, the sequence can be written as
\begin{equation}
    X =
    \begin{bmatrix}
        x_1\\
        x_2\\
        \vdots\\
        x_N
    \end{bmatrix}
    \in
    \mathbb{R}^{N\times d_{\mathrm{model}}},
    \qquad
    x_i\in\mathbb{R}^{1\times d_{\mathrm{model}}}.
    \label{eq:attn-input-matrix}
\end{equation}
Each row is the current representation of one token position. The word
``current'' is important: in a deep Transformer, a row of \(X\) is not
necessarily the original embedding of the token. It may already contain
information produced by preceding layers.

The representation required at position \(i\) generally depends on more than
the token at that position. A word may be disambiguated by another word in the
sentence; a pronoun may refer to an entity introduced earlier; and a question
may require information distributed across several parts of a document. A
model that transforms each token independently cannot directly exchange
information between positions through that transformation alone. Attention
addresses this requirement by allowing each position to retrieve a
content-dependent combination of information from other positions.

The mechanism separates two questions. The first is which positions are
relevant to the current position. The second is what information should be
transferred from those positions. Queries and keys determine the first
relationship, while values carry the information used in the second
operation. This separation is the basis of the query--key--value formulation
of attention described in the supplied Transformer survey
\parencite{ansar2024transformers}.

\begin{figure}[htbp]
    \centering
    % Figure source fragment: scaled dot-product attention data flow.
% This file is inserted inside a figure environment using \input.
\begin{center}
\setlength{\fboxsep}{5pt}
\renewcommand{\arraystretch}{1.2}
\begin{tabular}{c}
    \fbox{\parbox[c][0.8cm][c]{0.48\textwidth}{\centering
        \textbf{Input sequence}
        \(X\in\mathbb{R}^{N\times d_{\mathrm{model}}}\)}}\\[0.3em]
    \(\swarrow\hspace{1.0cm}\downarrow\hspace{1.0cm}\searrow\)\\[0.2em]
    \begin{tabular}{ccc}
        \fbox{\parbox[c][0.8cm][c]{0.21\textwidth}{\centering
            \textbf{Query}\\
            \(Q=XW_Q\)}} &
        \fbox{\parbox[c][0.8cm][c]{0.21\textwidth}{\centering
            \textbf{Key}\\
            \(K=XW_K\)}} &
        \fbox{\parbox[c][0.8cm][c]{0.21\textwidth}{\centering
            \textbf{Value}\\
            \(V=XW_V\)}}
    \end{tabular}\\[0.35em]
    \(\searrow\hspace{0.7cm}\swarrow\hspace{2.0cm}\downarrow\)\\[0.2em]
    \begin{tabular}{cc}
        \fbox{\parbox[c][0.85cm][c]{0.36\textwidth}{\centering
            \textbf{Scaled compatibility}\\
            \(S=QK^\top/\sqrt{d_k}\)}} &
        \fbox{\parbox[c][0.85cm][c]{0.27\textwidth}{\centering
            \textbf{Value matrix}\\
            \(V\)}}
    \end{tabular}\\[0.3em]
    \(\downarrow\hspace{2.7cm}\downarrow\)\\[0.2em]
    \begin{tabular}{cc}
        \fbox{\parbox[c][1.05cm][c]{0.36\textwidth}{\centering
            \textbf{Mask and row-wise softmax}\\
            \(P=\operatorname{softmax}_{\mathrm{row}}(S+M)\)}} &
        \fbox{\parbox[c][0.85cm][c]{0.27\textwidth}{\centering
            \textbf{Values retained}\\
            \(V\)}}
    \end{tabular}\\[0.3em]
    \(\searrow\hspace{0.8cm}\swarrow\)\\[0.2em]
    \fbox{\parbox[c][0.85cm][c]{0.56\textwidth}{\centering
        \textbf{Weighted value aggregation}\\
        \(O=PV\)}}
\end{tabular}
\end{center}

    \caption{Computational flow of scaled dot-product attention. The input
    representations are projected into queries, keys, and values; query--key
    compatibility scores are scaled and normalized; and the resulting
    attention weights determine the weighted aggregation of value vectors.
    The formulation follows the attention descriptions in
    \textcite{ansar2024transformers} and \textcite{dao2023flashattention2}.}
    \label{fig:attn-qkv-pipeline}
\end{figure}
% Alt text for accessibility register: A flow diagram showing an input
% sequence matrix branching into query, key, and value projections. Queries
% and transposed keys form a score matrix, which is scaled, masked if needed,
% and normalized by row-wise softmax. The resulting weights multiply values
% to produce contextualized output representations.

\subsection{Query, key, and value representations}
\label{subsec:attn-qkv-representations}

Attention begins by applying three learned linear projections to the input
matrix \(X\). For one attention head, define
\begin{align}
    Q &= XW_Q,
    & W_Q&\in\mathbb{R}^{d_{\mathrm{model}}\times d_k},
    \label{eq:attn-query-projection}\\
    K &= XW_K,
    & W_K&\in\mathbb{R}^{d_{\mathrm{model}}\times d_k},
    \label{eq:attn-key-projection}\\
    V &= XW_V,
    & W_V&\in\mathbb{R}^{d_{\mathrm{model}}\times d_v}.
    \label{eq:attn-value-projection}
\end{align}
The resulting matrices have dimensions
\begin{equation}
    Q,K\in\mathbb{R}^{N\times d_k},
    \qquad
    V\in\mathbb{R}^{N\times d_v}.
    \label{eq:attn-qkv-shapes}
\end{equation}
Here, \(d_k\) is the query and key dimension, and \(d_v\) is the value
dimension. The matrices \(W_Q\), \(W_K\), and \(W_V\) are trainable
parameters. The \(i\)-th rows of \(Q\), \(K\), and \(V\) are denoted by
\(q_i\), \(k_i\), and \(v_i\), respectively:
\begin{equation}
    q_i=x_iW_Q,\qquad
    k_i=x_iW_K,\qquad
    v_i=x_iW_V.
    \label{eq:attn-qkv-token}
\end{equation}

A query is the projected representation used by a position to express what
information it seeks. A key is the projected representation against which
queries are compared. A value is the projected content contributed by a
position when another position attends to it. These descriptions identify
their computational roles; they do not imply that queries, keys, or values
are literal symbolic questions, labels, or database records. They are learned
vectors whose geometry is shaped by optimization.

The query and key dimensions must agree because their compatibility is
computed by a dot product. The value dimension need not equal the key
dimension. This distinction allows the model to learn one representation
space for deciding where to retrieve information and another for representing
the information that is transferred.

\subsection{Scaled dot-product attention}
\label{subsec:attn-scaled-dot-product}

\paragraph{Compatibility scores.}
For a query \(q_i\) and key \(k_j\), the unnormalized compatibility score is
their inner product:
\begin{equation}
    s_{ij}=q_i k_j^\top
    =\sum_{r=1}^{d_k}q_{ir}k_{jr}.
    \label{eq:attn-dot-product-score}
\end{equation}
The dot product is a scalar measure of alignment in the learned query--key
space. It is not a probability and is not restricted to the interval
\([0,1]\). It may be positive, negative, or zero. For all query--key pairs,
the scores are computed by the matrix multiplication
\begin{equation}
    S=QK^\top,
    \qquad
    S\in\mathbb{R}^{N\times N}.
    \label{eq:attn-score-matrix}
\end{equation}
The element \(S_{ij}\) is the score assigned by query position \(i\) to key
position \(j\). Thus, dense self-attention constructs an \(N\times N\)
matrix of pairwise interactions.

\paragraph{Why the scores are scaled.}
The magnitude of an unscaled dot product depends on the dimension \(d_k\).
A statistical argument explains the scaling used in the standard
formulation. Assume, for this calculation, that the components of \(q_i\)
and \(k_j\) are mutually independent, have zero mean, and have unit
variance. Then
\begin{equation}
    \mathbb{E}[q_i k_j^\top]=0,
    \qquad
    \operatorname{Var}(q_i k_j^\top)=d_k.
    \label{eq:attn-dot-product-variance}
\end{equation}
The variance therefore increases linearly with the key dimension. Dividing
the dot product by \(\sqrt{d_k}\) gives a score with unit variance under
these assumptions:
\begin{equation}
    \operatorname{Var}
    \left(
        \frac{q_i k_j^\top}{\sqrt{d_k}}
    \right)
    =
    \frac{d_k}{d_k}=1.
    \label{eq:attn-scaled-variance}
\end{equation}
This is a scale-control argument under a simplifying statistical model, not
a claim that trained attention scores necessarily have unit variance.

The scaled score matrix is
\begin{equation}
    \widetilde S
    =
    \frac{QK^\top}{\sqrt{d_k}}.
    \label{eq:attn-scaled-score}
\end{equation}
The scaling reduces the tendency for score magnitudes to grow merely because
the projection dimension is larger. This is relevant because softmax can
become highly concentrated when its inputs have large differences. A
concentrated distribution can yield small derivatives for many score
components and can make optimization less well-conditioned. The scaled
dot-product formulation is the standard attention computation described in
the supplied surveys and in the FlashAttention-2 paper
\parencite{ansar2024transformers,dao2023flashattention2}.

\paragraph{Softmax normalization.}
The score matrix is converted into attention weights by applying softmax
independently to each row. For a vector
\(s=(s_1,\ldots,s_N)\), softmax is
\begin{equation}
    \operatorname{softmax}(s)_j
    =
    \frac{\exp(s_j)}
         {\displaystyle\sum_{r=1}^{N}\exp(s_r)}.
    \label{eq:attn-softmax-definition}
\end{equation}
The attention probability matrix is therefore
\begin{equation}
    P=
    \operatorname{softmax}_{\mathrm{row}}
    \left(
        \frac{QK^\top}{\sqrt{d_k}}
    \right),
    \qquad
    P\in\mathbb{R}^{N\times N}.
    \label{eq:attn-probability-matrix}
\end{equation}
For every row \(i\),
\begin{equation}
    P_{ij}\geq 0,
    \qquad
    \sum_{j=1}^{N}P_{ij}=1.
    \label{eq:attn-probability-properties}
\end{equation}
Thus, each row is a probability distribution over the key positions
available to the corresponding query. The weights are relative: changing
one score changes the normalized weights of the other positions in that row
as well.

A simple numerical calculation illustrates the normalization. If a query
has two scaled scores \(s_1=0\) and \(s_2=\ln 2\), then
\begin{equation}
    \operatorname{softmax}(0,\ln 2)
    =
    \left(
        \frac{1}{3},\frac{2}{3}
    \right).
    \label{eq:attn-softmax-example}
\end{equation}
The second key receives twice the normalized weight of the first key. This
does not mean that its value vector is twice as large; it means that its
contribution to the weighted aggregation is multiplied by twice the
coefficient.

For numerical stability, softmax is commonly evaluated after subtracting
the largest score in the row. If
\(m=\max_j s_j\), then
\begin{equation}
    \frac{\exp(s_j)}
         {\sum_r\exp(s_r)}
    =
    \frac{\exp(s_j-m)}
         {\sum_r\exp(s_r-m)}.
    \label{eq:attn-stable-softmax}
\end{equation}
The equality follows because the same factor \(\exp(-m)\) appears in the
numerator and denominator and cancels. The transformed exponents are
non-positive, which reduces the risk of overflow. This identity is also
the basis of the online softmax updates used by FlashAttention
\parencite{dao2023flashattention2}.

\paragraph{Weighted aggregation of values.}
The attention weights determine how value vectors are combined:
\begin{equation}
    O=PV,
    \qquad
    O\in\mathbb{R}^{N\times d_v}.
    \label{eq:attn-value-aggregation}
\end{equation}
For position \(i\), the output is
\begin{equation}
    o_i=\sum_{j=1}^{N}P_{ij}v_j.
    \label{eq:attn-token-aggregation}
\end{equation}
Because the weights are non-negative and sum to one, this output is a
convex combination of the value vectors included in that row's
normalization. It therefore lies in their convex hull. The output is
context-dependent because the weights depend on the current query and all
available keys.

Combining the score, normalization, and aggregation stages gives the
scaled dot-product attention operation:
\begin{equation}
    \boxed{
    \operatorname{Attention}(Q,K,V)
    =
    \operatorname{softmax}_{\mathrm{row}}
    \left(
        \frac{QK^\top}{\sqrt{d_k}}
    \right)V
    }.
    \label{eq:attn-scaled-dot-product}
\end{equation}
The equation separates the decision about which positions to retrieve from
the transformation of the information being retrieved.

\subsection{Self-attention, cross-attention, and causal masking}
\label{subsec:attn-self-cross-causal}

\paragraph{Self-attention.}
Self-attention is the configuration in which queries, keys, and values are
all projected from the same input sequence:
\begin{equation}
    Q=XW_Q,\qquad
    K=XW_K,\qquad
    V=XW_V.
    \label{eq:attn-self-qkv}
\end{equation}
The output is
\begin{equation}
    \operatorname{SelfAttention}(X)
    =
    \operatorname{softmax}_{\mathrm{row}}
    \left(
        \frac{XW_QW_K^\top X^\top}{\sqrt{d_k}}
    \right)XW_V.
    \label{eq:attn-self-attention}
\end{equation}
Every position can use its query to identify relevant positions in the same
sequence. The interaction pattern is content-dependent because it is
calculated from the current learned representations, rather than fixed
solely by token distance. Self-attention is a central component of
Transformer-based language models and is described as such in the supplied
Transformer survey \parencite{ansar2024transformers}.

\paragraph{Cross-attention.}
Queries do not have to originate from the same sequence as keys and values.
Let \(X_q\) be a query sequence of length \(N_q\), and let \(X_c\) be a
context sequence of length \(N_c\). Cross-attention uses
\begin{align}
    Q&=X_qW_Q,
    & Q&\in\mathbb{R}^{N_q\times d_k},
    \label{eq:attn-cross-query}\\
    K&=X_cW_K,
    & K&\in\mathbb{R}^{N_c\times d_k},
    \label{eq:attn-cross-key}\\
    V&=X_cW_V,
    & V&\in\mathbb{R}^{N_c\times d_v}.
    \label{eq:attn-cross-value}
\end{align}
The score matrix then has shape
\begin{equation}
    S=QK^\top
    \in\mathbb{R}^{N_q\times N_c},
    \label{eq:attn-cross-score-shape}
\end{equation}
and the output is
\begin{equation}
    O=
    \operatorname{softmax}_{\mathrm{row}}
    \left(
        \frac{QK^\top}{\sqrt{d_k}}
    \right)V
    \in\mathbb{R}^{N_q\times d_v}.
    \label{eq:attn-cross-output}
\end{equation}
Each query position can therefore retrieve information from the context
sequence even when the two sequences have different lengths. This is the
form of attention used when a decoder retrieves information from encoder
representations in an encoder--decoder arrangement. The supplied GQA study
also evaluates multi-query and grouped-query variants in decoder
cross-attention as well as decoder self-attention
\parencite{ainslie2023gqa}.

\paragraph{Causal masking.}
Autoregressive language modeling imposes an information constraint:
the prediction at position \(i\) must not depend on tokens at positions
greater than \(i\). Without a mask, ordinary self-attention could use
future tokens during training, producing information leakage relative to
the next-token prediction task.

Define the additive causal mask \(M\in\mathbb{R}^{N\times N}\) by
\begin{equation}
    M_{ij}
    =
    \begin{cases}
        0, & j\leq i,\\
        -\infty, & j>i.
    \end{cases}
    \label{eq:attn-causal-mask}
\end{equation}
The masked score matrix is
\begin{equation}
    S_{\mathrm{causal}}
    =
    \frac{QK^\top}{\sqrt{d_k}}+M.
    \label{eq:attn-causal-scores}
\end{equation}
Because \(\exp(-\infty)=0\), future positions receive zero attention
probability. The row-\(i\) distribution is normalized only over positions
\(1,\ldots,i\):
\begin{equation}
    P_{ij}
    =
    \begin{cases}
        \displaystyle
        \frac{\exp(\widetilde S_{ij})}
             {\sum_{r=1}^{i}\exp(\widetilde S_{ir})},
        & j\leq i,\\[10pt]
        0, & j>i,
    \end{cases}
    \label{eq:attn-causal-probabilities}
\end{equation}
where \(\widetilde S=QK^\top/\sqrt{d_k}\). Causal attention is consequently
lower triangular. The first position can use only itself; the second can
use the first and second positions; and position \(i\) can use the prefix
ending at \(i\).

The mask defines the permitted information flow. During autoregressive
generation, the model processes the available prefix, predicts a next
token, appends that token, and makes its representation available to
subsequent positions. The causal restriction is therefore an architectural
condition for autoregressive prediction, not merely a numerical
optimization. Implementations can exploit the triangular structure by
skipping blocks that lie entirely in the masked region, as discussed in
the FlashAttention-2 paper \parencite{dao2023flashattention2}.

\subsection{Multi-head attention}
\label{subsec:attn-multi-head}

A single attention head uses one learned query--key compatibility space and
one value space. A language sequence, however, can contain several
simultaneous kinds of relationships. A token may need to retrieve a nearby
syntactic dependency, an entity introduced much earlier, or information
associated with a task-specific pattern. Requiring all such relationships
to be represented through one projection space may constrain the range of
interactions the layer can learn.

Multi-head attention addresses this limitation by performing several
attention computations in parallel. Let \(H\) be the number of heads. For
head \(h\), define
\begin{align}
    Q_h&=XW_Q^{(h)},&
    K_h&=XW_K^{(h)},&
    V_h&=XW_V^{(h)},
    \label{eq:attn-head-qkv}\\
    \operatorname{head}_h
    &=
    \operatorname{softmax}_{\mathrm{row}}
    \left(
        \frac{Q_hK_h^\top}{\sqrt{d_k}}+M
    \right)V_h.
    \label{eq:attn-head-output}
\end{align}
Each head has its own learned projection matrices and can therefore learn
a different compatibility function and value transformation.

In the conventional equal-dimension construction,
\begin{equation}
    d_k=d_v=d_h=\frac{d_{\mathrm{model}}}{H}.
    \label{eq:attn-head-dimension}
\end{equation}
Each head output has shape
\begin{equation}
    \operatorname{head}_h
    \in\mathbb{R}^{N\times d_h}.
    \label{eq:attn-head-output-shape}
\end{equation}
The outputs are concatenated along their feature dimensions:
\begin{equation}
    C=
    \operatorname{Concat}
    \left(
        \operatorname{head}_1,\ldots,\operatorname{head}_H
    \right)
    \in
    \mathbb{R}^{N\times Hd_h}
    =
    \mathbb{R}^{N\times d_{\mathrm{model}}}.
    \label{eq:attn-head-concatenation}
\end{equation}
A learned output projection maps the concatenated result into the model
representation space:
\begin{equation}
    \boxed{
    \operatorname{MHA}(X)=CW_O,
    \qquad
    W_O\in
    \mathbb{R}^{d_{\mathrm{model}}\times d_{\mathrm{model}}}.
    }
    \label{eq:attn-mha-output}
\end{equation}
The output projection allows information from different head feature
dimensions to be combined. The multi-head construction is summarized in
\cref{fig:attn-multi-head}; the supplied Transformer survey describes the
same parallel-head, concatenation, and output-projection structure
\parencite{ansar2024transformers}.

\begin{figure}[htbp]
    \centering
    % Figure source fragment: multi-head attention.
% This file is inserted inside a figure environment using \input.
\begin{center}
\setlength{\fboxsep}{5pt}
\renewcommand{\arraystretch}{1.2}
\begin{tabular}{c}
    \fbox{\parbox[c][0.75cm][c]{0.55\textwidth}{\centering
        \textbf{Shared input sequence} \(X\)}}\\[0.3em]
    \(\swarrow\hspace{0.8cm}\downarrow\hspace{0.8cm}\searrow\)\\[0.2em]
    \begin{tabular}{ccccc}
        \fbox{\parbox[c][0.95cm][c]{0.16\textwidth}{\centering
            \textbf{Head 1}\\
            \(Q_h,K_h,V_h\)}} &
        \fbox{\parbox[c][0.95cm][c]{0.16\textwidth}{\centering
            \textbf{Head 2}\\
            \(Q_h,K_h,V_h\)}} &
        \(\cdots\) &
        \fbox{\parbox[c][0.95cm][c]{0.16\textwidth}{\centering
            \textbf{Head \(H-1\)}\\
            \(Q_h,K_h,V_h\)}} &
        \fbox{\parbox[c][0.95cm][c]{0.16\textwidth}{\centering
            \textbf{Head \(H\)}\\
            \(Q_h,K_h,V_h\)}}
    \end{tabular}\\[0.35em]
    \(\downarrow\hspace{0.45cm}\downarrow\hspace{0.65cm}
      \cdots\hspace{0.65cm}\downarrow\hspace{0.45cm}\downarrow\)\\[0.2em]
    \fbox{\parbox[c][0.85cm][c]{0.67\textwidth}{\centering
        \textbf{Concatenation}\\
        \(C=\operatorname{Concat}(\operatorname{head}_1,\ldots,
        \operatorname{head}_H)\)}}\\[0.3em]
    \(\downarrow\;W_O\)\\[0.2em]
    \fbox{\parbox[c][0.85cm][c]{0.55\textwidth}{\centering
        \textbf{Output projection}\\
        \(\operatorname{MHA}(X)=CW_O\)}}
\end{tabular}
\end{center}

    \caption{Multi-head attention. The same input sequence is projected into
    separate query, key, and value representations for each head. The head
    outputs are concatenated and passed through a learned output projection.
    Based on the multi-head formulation in
    \textcite{ansar2024transformers}.}
    \label{fig:attn-multi-head}
\end{figure}
% Alt text for accessibility register: A diagram showing one input matrix
% feeding several parallel attention heads. Each head produces an output;
% these outputs are concatenated and passed through one output projection
% to form the multi-head attention result.

\paragraph{Parameter count.}
Assume that all heads have dimension \(d_h=d_{\mathrm{model}}/H\), and
ignore projection biases. Across all heads, the query projection weights
contain
\begin{equation}
    P_Q=H d_{\mathrm{model}}d_h
    =d_{\mathrm{model}}^2
    \label{eq:attn-mha-query-parameters}
\end{equation}
parameters. The key and value projections have the same count:
\begin{equation}
    P_K=P_V=d_{\mathrm{model}}^2.
    \label{eq:attn-mha-key-value-parameters}
\end{equation}
The output projection contributes another \(d_{\mathrm{model}}^2\)
parameters. Thus, under these assumptions,
\begin{equation}
    P_{\mathrm{MHA}}
    =
    P_Q+P_K+P_V+P_O
    =
    4d_{\mathrm{model}}^2.
    \label{eq:attn-mha-parameter-count}
\end{equation}
This count is independent of \(H\) when the total model dimension remains
fixed and the per-head dimension is reduced proportionally. Increasing
the number of heads therefore changes the partition of the projection
space, but does not by itself increase this idealized total projection
parameter count.

\paragraph{Arithmetic cost.}
For one head, computing \(Q_hK_h^\top\) requires approximately
\(2N^2d_h\) floating-point operations when a multiplication and an
addition are counted separately. Computing \(P_hV_h\) requires another
\(2N^2d_h\). Summing over \(H\) heads gives
\begin{equation}
    C_{\mathrm{attention}}
    \approx
    4HN^2d_h
    =
    4N^2d_{\mathrm{model}}
    \label{eq:attn-mha-flops}
\end{equation}
for the two principal matrix multiplications of one sequence. This
expression excludes the query, key, value, and output projections, the
softmax, masking, and other layer operations. Its principal implication
is the quadratic dependence on \(N\).

\subsection{Backpropagation through attention}
\label{subsec:attn-backpropagation}

Attention is a differentiable computation whose projection matrices are
learned through the loss function. The backward pass determines how the
loss changes when the query, key, and value representations change. This
derivation also explains why storing or recomputing attention probabilities
is important for memory-efficient training.

Consider the forward equations
\begin{equation}
    S=\frac{QK^\top}{\sqrt{d_k}},
    \qquad
    P=\operatorname{softmax}_{\mathrm{row}}(S),
    \qquad
    O=PV.
    \label{eq:attn-backward-forward}
\end{equation}
Let \(\mathcal{L}\) be a scalar loss and define the upstream gradient
\begin{equation}
    G=\frac{\partial\mathcal{L}}{\partial O}
    \in\mathbb{R}^{N\times d_v}.
    \label{eq:attn-upstream-gradient}
\end{equation}

The value gradient follows directly from \(O=PV\):
\begin{equation}
    \frac{\partial\mathcal{L}}{\partial V}
    =
    P^\top G.
    \label{eq:attn-gradient-value}
\end{equation}
The probability gradient is
\begin{equation}
    \frac{\partial\mathcal{L}}{\partial P}
    =
    GV^\top
    \in\mathbb{R}^{N\times N}.
    \label{eq:attn-gradient-probability}
\end{equation}

The softmax derivative must account for the coupling introduced by
normalization. For one row \(p=\operatorname{softmax}(s)\), the Jacobian
is
\begin{equation}
    J_{\mathrm{softmax}}(p)
    =
    \operatorname{diag}(p)-pp^\top.
    \label{eq:attn-softmax-jacobian}
\end{equation}
This follows from
\begin{equation}
    \frac{\partial p_j}{\partial s_k}
    =
    p_j(\delta_{jk}-p_k),
    \label{eq:attn-softmax-component-derivative}
\end{equation}
where \(\delta_{jk}\) is the Kronecker delta. If
\(D_P=\partial\mathcal{L}/\partial P\), the score gradient for each row is
\begin{equation}
    d s
    =
    \left(
        \operatorname{diag}(p)-pp^\top
    \right)d p.
    \label{eq:attn-softmax-backward}
\end{equation}
Equivalently, in component form,
\begin{equation}
    \frac{\partial\mathcal{L}}{\partial s_j}
    =
    p_j
    \left(
        \frac{\partial\mathcal{L}}{\partial p_j}
        -
        \sum_{r=1}^{N}
        p_r\frac{\partial\mathcal{L}}{\partial p_r}
    \right).
    \label{eq:attn-softmax-backward-component}
\end{equation}
The gradient of one score therefore depends on the probability and
upstream gradient of every score in the same row.

Let
\(D=\partial\mathcal{L}/\partial S\) denote the matrix formed by applying
this row-wise derivative. The gradients of the scaled query--key product
are
\begin{align}
    \frac{\partial\mathcal{L}}{\partial Q}
    &=
    \frac{DK}{\sqrt{d_k}},
    \label{eq:attn-gradient-query}\\
    \frac{\partial\mathcal{L}}{\partial K}
    &=
    \frac{D^\top Q}{\sqrt{d_k}}.
    \label{eq:attn-gradient-key}
\end{align}
For a masked score, the corresponding probability is zero and its
gradient contribution is zero. The complete backward pass is consequently
\begin{align}
    dV &= P^\top G,
    \label{eq:attn-backward-dv}\\
    dP &= GV^\top,
    \label{eq:attn-backward-dp}\\
    dS &= \operatorname{SoftmaxBackward}(P,dP),
    \label{eq:attn-backward-ds}\\
    dQ &= \frac{dSK}{\sqrt{d_k}},
    \label{eq:attn-backward-dq}\\
    dK &= \frac{dS^\top Q}{\sqrt{d_k}}.
    \label{eq:attn-backward-dk}
\end{align}

The gradients then propagate to the learned projection matrices. Since
\(Q=XW_Q\), \(K=XW_K\), and \(V=XW_V\), the chain rule gives
\begin{align}
    \frac{\partial\mathcal{L}}{\partial W_Q}
    &=
    X^\top\frac{\partial\mathcal{L}}{\partial Q},
    \label{eq:attn-gradient-wq}\\
    \frac{\partial\mathcal{L}}{\partial W_K}
    &=
    X^\top\frac{\partial\mathcal{L}}{\partial K},
    \label{eq:attn-gradient-wk}\\
    \frac{\partial\mathcal{L}}{\partial W_V}
    &=
    X^\top\frac{\partial\mathcal{L}}{\partial V}.
    \label{eq:attn-gradient-wv}
\end{align}
The contribution to the input gradient through these three projections is
\begin{equation}
    \frac{\partial\mathcal{L}}{\partial X}
    =
    \frac{\partial\mathcal{L}}{\partial Q}W_Q^\top
    +
    \frac{\partial\mathcal{L}}{\partial K}W_K^\top
    +
    \frac{\partial\mathcal{L}}{\partial V}W_V^\top.
    \label{eq:attn-gradient-input}
\end{equation}
For multi-head attention, the output projection combines the
head-specific contributions before they return to the shared input
representation. Training therefore adjusts both the information-selection
spaces and the value transformations. The same backward dependencies
motivate the recomputation strategy used by FlashAttention
\parencite{dao2023flashattention2}.

\subsection{Attention-head interpretation and functional specialization}
\label{subsec:attn-head-interpretability}

The multi-head formulation explains how several attention computations
can operate in parallel, but it does not establish what a particular head
has learned to do. Investigating that question belongs to mechanistic
interpretability: the study of internal computations through parameters,
activations, intermediate representations, and controlled interventions.
The supplied survey by \textcite{zheng2024attentionheads} organizes
reported attention-head functions through a four-stage framework:
Knowledge Recalling, In-Context Identification, Latent Reasoning, and
Expression Preparation. The framework is a way to organize experimental
findings; it is not evidence that a Transformer contains four explicit
modules corresponding to human cognitive stages.

\paragraph{The QK and OV transformations.}
For one head \(h\), let \(W_Q^{(h)}\), \(W_K^{(h)}\), and \(W_V^{(h)}\)
be its projection matrices. Let \(W_O^{(h)}\) denote the slice of the
multi-head output projection that maps this head's output into the model
dimension. The head's contribution can be expressed as
\begin{equation}
    \operatorname{Head}_h(X)
    =
    \operatorname{softmax}
    \left(
        \frac{
            XW_Q^{(h)}W_K^{(h)\top}X^\top
        }{\sqrt{d_h}}+M
    \right)
    XW_V^{(h)}W_O^{(h)}.
    \label{eq:attn-head-expanded}
\end{equation}
This expression exposes two composite transformations. The query--key
matrix is
\begin{equation}
    W_{QK}^{(h)}
    =
    W_Q^{(h)}W_K^{(h)\top},
    \label{eq:attn-qk-matrix}
\end{equation}
which determines how token representations are compared to produce
attention scores. The value--output matrix is
\begin{equation}
    W_{OV}^{(h)}
    =
    W_V^{(h)}W_O^{(h)},
    \label{eq:attn-ov-matrix}
\end{equation}
which determines how attended information is transformed before it is
written into the model representation space. The QK transformation is
therefore associated with information selection, while the OV
transformation is associated with information transfer. The distinction
is central to the interpretability framework in
\textcite{zheng2024attentionheads}.

\paragraph{Residual-stream contribution.}
An attention sublayer usually adds an update to the representation
carried through the Transformer by residual connections. Omitting
architecture-dependent normalization details, this update can be written as
\begin{equation}
    R_{\ell}^{\mathrm{after}}
    =
    R_{\ell}^{\mathrm{before}}
    +
    \sum_{h=1}^{H}\operatorname{Head}_{\ell,h}(R_{\ell}^{\mathrm{before}}).
    \label{eq:attn-residual-stream}
\end{equation}
The residual stream is the shared representation pathway through which
information from different heads and layers can accumulate. A head reads
the representation available in this stream, performs its query--key
selection and value transformation, and writes an update back into the
stream. Subsequent heads and layers can then use the modified
representation. The exact placement of normalization and the ordering of
sublayers depend on the Transformer variant; the equation isolates the
residual contribution rather than specifying one universal block layout.

\paragraph{Knowledge Recalling.}
The Knowledge Recalling stage concerns retrieval of information associated
with the model's parametric knowledge. The survey identifies the Memory
Head as an example: it can retrieve attributes associated with entities
present in the input and write the retrieved information into the residual
stream. In task-specific settings, the survey also describes Constant
Heads, Single Letter Heads, and Negative Heads in multiple-choice or
binary-decision contexts. These examples illustrate that some heads are
associated with retrieval from learned parameters, while others are
associated with the structure of a particular task
\parencite{zheng2024attentionheads}.

\paragraph{In-Context Identification.}
The In-Context Identification stage concerns identifying structural,
syntactic, and semantic information in the current context. The survey
discusses Previous Heads and Positional Heads as examples that track
positional relationships, Duplicate Heads as heads that respond to
repeated content, and Retrieval Heads as heads that locate specific
information in longer contexts. Syntactic Heads are reported to identify
relationships such as subjects, objects, and modifiers. Subword Merge
Heads are associated with integrating subword representations, while
Name Mover and Letter Mover Heads transfer entity or letter information
toward positions used in subsequent computation. Context Heads and
Content Gatherer Heads identify task-relevant content; Sentiment
Summarizers and Semantic Induction Heads are associated with sentiment
and semantic relationships. These are reported functional descriptions
from the surveyed studies, not universal properties of every model's
attention heads \parencite{zheng2024attentionheads}.

\paragraph{Latent Reasoning.}
The Latent Reasoning stage concerns the integration of information
identified in the context and the use of learned patterns to produce
intermediate conclusions. Induction Heads are a representative example.
The survey describes a pattern of the form
\begin{equation}
    \ldots [A][B]\ldots[A],
    \label{eq:attn-induction-pattern}
\end{equation}
where the first occurrence of \([A]\) is followed by \([B]\). An Induction
Head can use information about the earlier continuation to support the
prediction of \([B]\) after the later occurrence of \([A]\). The survey
also discusses Summary Readers and Function Vectors in task recognition,
as well as heads associated with truthfulness, accuracy, consistency,
iteration, successor operations, and inhibition of selected signals.
These examples show that head functions can involve both the retrieval
and combination of information and the suppression of competing
information \parencite{zheng2024attentionheads}.

\paragraph{Expression Preparation.}
The Expression Preparation stage concerns preparing information for
expression in the model's output. The survey describes Mixed Heads as
aggregating information passed from other heads, Amplification and
Correct Heads as strengthening signals associated with a selected answer,
and Coherence and Faithfulness Heads as being associated with linguistic
consistency and alignment between reasoning-related information and
instructions. The framework also allows that some tasks may not require
a distinct expression-preparation operation: information already present
in the residual stream may be sufficient to determine the next output
token \parencite{zheng2024attentionheads}.

The four stages can interact and need not occur as a rigid one-way
sequence. A model may return to contextual identification after retrieving
knowledge, or retrieve additional information during reasoning. The
framework is therefore best interpreted as a descriptive organization of
reported head functions rather than a literal computational pipeline.
\cref{fig:attn-head-framework} summarizes the four categories and
illustrative functions discussed in the survey.

\begin{figure}[htbp]
    \centering
    % Figure source fragment: four-stage attention-head interpretability framework.
% This file is inserted inside a figure environment using \input.
\begin{center}
\setlength{\fboxsep}{5pt}
\renewcommand{\arraystretch}{1.2}
\begin{tabular}{cc}
    \fbox{\parbox[t][1.65cm][t]{0.42\textwidth}{\centering
        \textbf{Knowledge Recalling}\\[0.2em]
        \footnotesize Memory Head; Constant Head;\\
        Single Letter Head; Negative Head.}} &
    \fbox{\parbox[t][1.65cm][t]{0.42\textwidth}{\centering
        \textbf{In-Context Identification}\\[0.2em]
        \footnotesize Previous Head; Retrieval Head;\\
        Syntactic Head; Name Mover Head.}}\\[0.35em]
    \fbox{\parbox[t][1.65cm][t]{0.42\textwidth}{\centering
        \textbf{Latent Reasoning}\\[0.2em]
        \footnotesize Induction Head; Function Vector;\\
        Truthfulness Head; Successor Head.}} &
    \fbox{\parbox[t][1.65cm][t]{0.42\textwidth}{\centering
        \textbf{Expression Preparation}\\[0.2em]
        \footnotesize Mixed Head; Amplification Head;\\
        Correct Head; Coherence and Faithfulness Heads.}}
\end{tabular}

\par\smallskip
\footnotesize
The stages are descriptive categories and can interact; they do not form
a rigid computational pipeline.
\normalsize
\end{center}

    \caption{Four-stage organization of reported attention-head functions:
    Knowledge Recalling, In-Context Identification, Latent Reasoning, and
    Expression Preparation. The categories are an interpretability
    framework and may interact rather than forming a rigid sequential
    pipeline. Based on \textcite{zheng2024attentionheads}.}
    \label{fig:attn-head-framework}
\end{figure}
% Alt text for accessibility register: Four connected boxes show Knowledge
% Recalling, In-Context Identification, Latent Reasoning, and Expression
% Preparation. Each box contains representative attention-head functions;
% arrows indicate that information can move between stages.

\paragraph{Experimental investigation.}
The survey describes experiments that test whether a proposed head
function contributes to model behavior. Activation patching replaces
selected activations with values from another input, a baseline, or a
perturbed run, and measures the effect on the output. Ablation instead removes or suppresses a component and
observes the resulting behavioral change. The two methods are related but
not identical: patching substitutes an activation, whereas ablation removes
a component's contribution. The survey also describes the logit lens,
which uses the unembedding transformation to map intermediate
representations to vocabulary logits and thereby examine how an
intermediate state relates to candidate output tokens
\parencite{zheng2024attentionheads}.

These interventions provide evidence about the role of a component under
the tested conditions. They do not, by themselves, establish that a
proposed function generalizes to every task, model family, or context.
The survey identifies generalizability, transferability, interactions
among multiple heads, and the need for stronger theoretical accounts as
important limitations in the study of attention-head functions
\parencite{zheng2024attentionheads}.

\subsection{Computational and memory complexity of dense attention}
\label{subsec:attn-computational-memory-complexity}

The dimensions of the score and probability matrices explain the principal
scaling limitation of dense attention. For one head, the two main matrix
multiplications are
\begin{equation}
    QK^\top:
    (N\times d_k)(d_k\times N)
    \longrightarrow
    N\times N
    \label{eq:attn-complexity-score-shape}
\end{equation}
and
\begin{equation}
    PV:
    (N\times N)(N\times d_v)
    \longrightarrow
    N\times d_v.
    \label{eq:attn-complexity-value-shape}
\end{equation}
Their combined arithmetic cost is
\begin{equation}
    \mathcal{O}(N^2d_k+N^2d_v).
    \label{eq:attn-complexity-one-head}
\end{equation}
When \(d_k=d_v=d_h\), this is
\begin{equation}
    \mathcal{O}(N^2d_h)
    \label{eq:attn-complexity-one-head-equal}
\end{equation}
up to constant factors. Across \(H\) heads with
\(Hd_h=d_{\mathrm{model}}\), the leading sequence-length dependence is
\begin{equation}
    \boxed{
    \mathcal{O}(N^2d_{\mathrm{model}})
    }.
    \label{eq:attn-complexity-mha}
\end{equation}

The quadratic dependence can be illustrated numerically. If the sequence
length increases from \(N=2{,}048\) to \(N=8{,}192\), the ratio of the
numbers of pairwise positions is
\begin{equation}
    \frac{8{,}192^2}{2{,}048^2}
    =
    4^2
    =
    16.
    \label{eq:attn-quadratic-example}
\end{equation}
Thus, at fixed head and model dimensions, the principal dense-attention
matrix multiplications require sixteen times as much arithmetic after a
fourfold increase in sequence length.

\paragraph{Intermediate memory.}
The score matrix \(S\) and probability matrix \(P\) each contain \(N^2\)
elements per head and sequence. If a conventional implementation
materializes both matrices for a batch of \(B\) sequences and \(H\) heads,
the storage for these two intermediates is
\begin{equation}
    M_{\mathrm{intermediate}}
    =
    2BHN^2b
    \quad\text{bytes},
    \label{eq:attn-intermediate-memory}
\end{equation}
where \(b\) is the number of bytes per stored element. The precise peak
allocation depends on the implementation and which intermediates are
retained for the backward pass, but the quadratic dependence on \(N\)
remains the central difficulty. The FlashAttention-2 paper identifies
materialization of the score and probability matrices in high-bandwidth
memory as a major source of memory use and data movement
\parencite{dao2023flashattention2}.

\paragraph{The key--value cache during autoregressive inference.}
Training processes multiple positions in parallel subject to the causal
mask. Autoregressive inference instead generates one token at a time.
At each decoding step, the new query must interact with the keys and
values of preceding positions. Recomputing all preceding key and value
projections at every step would repeat work, so decoder implementations
commonly retain the previously computed keys and values in a
key--value cache, or KV cache.

Let \(L\) be the number of Transformer layers, \(B\) the batch size,
\(N\) the number of cached positions, \(H_{\mathrm{KV}}\) the number of
key--value heads, \(d_h\) the dimension of each such head, and \(b\) the
number of bytes per element. The cache contains keys and values, so its
memory requirement is
\begin{equation}
    \boxed{
    M_{\mathrm{KV}}
    =
    2LBNH_{\mathrm{KV}}d_hb
    }.
    \label{eq:attn-kv-cache-memory}
\end{equation}
The cache grows linearly with sequence length, but its total size can be
large when the model has many layers, a large batch, or many key--value
heads. During decoding, the keys and values must also be transferred from
memory for the current query computation. Consequently, inference speed
can be constrained by memory bandwidth, not only by arithmetic throughput.
This is the motivation for reducing the number of key--value heads in
multi-query and grouped-query attention
\parencite{ainslie2023gqa}.

\subsection{Multi-query and grouped-query attention}
\label{subsec:attn-mqa-gqa}

\paragraph{Multi-query attention.}
Conventional multi-head attention uses a separate key and value head for
each query head. Multi-query attention (MQA) retains multiple query heads
but shares one key head and one value head across all of them. If there
are \(H\) query heads, then
\begin{equation}
    H_Q=H,
    \qquad
    H_K=H_V=1.
    \label{eq:attn-mqa-head-counts}
\end{equation}
The output of query head \(h\) is
\begin{equation}
    \operatorname{head}_h
    =
    \operatorname{softmax}
    \left(
        \frac{Q_hK^\top}{\sqrt{d_h}}+M
    \right)V.
    \label{eq:attn-mqa-head}
\end{equation}
The query projection remains head-specific, but the same key and value
representations are used by all query heads.

For conventional MHA, the KV cache contains
\(2LBNHd_hb\) bytes. For MQA, it contains
\(2LBNd_hb\) bytes. Under the same dimensions, the reduction factor is
\begin{equation}
    R_{\mathrm{MQA}}
    =
    \frac{2LBNHd_hb}{2LBNd_hb}
    =
    H.
    \label{eq:attn-mqa-cache-reduction}
\end{equation}
This reduction decreases the amount of key--value data that must be stored
and loaded during autoregressive inference. It does not remove the
query-head computations. Sharing key and value representations can,
however, reduce representational flexibility and may cause quality
degradation or training instability. These trade-offs motivate an
intermediate design \parencite{ainslie2023gqa}.

\paragraph{Grouped-query attention.}
Grouped-query attention (GQA) retains \(H\) query heads but partitions
them into \(G\) groups, each sharing one key head and one value head.
Assume \(G\) divides \(H\). Each group then contains \(H/G\) query heads:
\begin{equation}
    H_Q=H,
    \qquad
    H_K=H_V=G.
    \label{eq:attn-gqa-head-counts}
\end{equation}
Let \(g(h)\) identify the key--value group assigned to query head \(h\).
The head output is
\begin{equation}
    \operatorname{head}_h
    =
    \operatorname{softmax}
    \left(
        \frac{
            Q_hK_{g(h)}^\top
        }{\sqrt{d_h}}+M
    \right)V_{g(h)}.
    \label{eq:attn-gqa-head}
\end{equation}
The KV-cache memory becomes
\begin{equation}
    M_{\mathrm{KV,GQA}}
    =
    2LBNGd_hb.
    \label{eq:attn-gqa-cache}
\end{equation}
Relative to MHA, the cache reduction factor is
\begin{equation}
    R_{\mathrm{GQA}}
    =
    \frac{H}{G}.
    \label{eq:attn-gqa-reduction}
\end{equation}
The limiting cases recover the two familiar architectures:
\begin{equation}
    G=1
    \Longrightarrow
    \operatorname{GQA}=\operatorname{MQA},
    \qquad
    G=H
    \Longrightarrow
    \operatorname{GQA}=\operatorname{MHA}.
    \label{eq:attn-gqa-boundaries}
\end{equation}
Intermediate values of \(G\) trade key--value memory against the number
of independently represented key--value groups.

For example, if a model has \(H=32\) query heads and uses \(G=8\)
key--value groups, then
\begin{equation}
    R_{\mathrm{GQA}}=\frac{32}{8}=4.
    \label{eq:attn-gqa-numerical-example}
\end{equation}
The KV cache is one quarter the size of the corresponding MHA cache,
while remaining eight times the size of an MQA cache under the same
head-dimension assumptions. The query-head count remains 32.

\paragraph{Converting and uptraining existing checkpoints.}
\textcite{ainslie2023gqa} propose converting a multi-head checkpoint into
an MQA or GQA checkpoint and then performing additional pre-training so
that the model can adapt to the changed key--value structure. For GQA,
the original key and value projection matrices are mean-pooled within
each new group. If \(\mathcal{G}_g\) is the set of original heads assigned
to group \(g\), the initialization is
\begin{align}
    W_K^{(g)}
    &=
    \frac{G}{H}
    \sum_{h\in\mathcal{G}_g}W_K^{(h)},
    \label{eq:attn-gqa-key-mean}\\
    W_V^{(g)}
    &=
    \frac{G}{H}
    \sum_{h\in\mathcal{G}_g}W_V^{(h)}.
    \label{eq:attn-gqa-value-mean}
\end{align}
The group mean retains contributions from the original heads rather than
selecting one head or randomly initializing the new shared projection.
The converted checkpoint is then uptrained using a fraction of the
original pre-training procedure. The authors report experiments with an
uptraining proportion of five percent of the original pre-training steps.

Their T5-XXL experiments report an inference time of \(1.51\) ms per
sample and an average task score of \(47.2\) for MHA-XXL. The uptrained
MQA-XXL model reports \(0.24\) ms and \(46.6\), while GQA-8-XXL reports
\(0.28\) ms and \(47.1\). These results illustrate the reported trade-off:
GQA-8 approaches the MHA-XXL average score while retaining inference
time close to MQA. The values are specific to the models, tasks, and
hardware configuration used in the study; they are not universal
performance guarantees \parencite{ainslie2023gqa}.

\begin{table}[htbp]
    \centering
    \caption{Selected T5-XXL results reported for MHA, MQA, and GQA.
    Inference time is the reported average time per sample; average
    performance is the reported average over the evaluated tasks.
    Source: \textcite{ainslie2023gqa}, Table 1.}
    \label{tab:attn-gqa-results}
    \small
    \begin{tabularx}{\textwidth}{l c c X}
        \toprule
        Configuration
        & Inference time (ms)
        & Average performance
        & Key--value structure \\
        \midrule
        MHA-XXL
        & 1.51
        & 47.2
        & One key--value head per query head \\
        MQA-XXL
        & 0.24
        & 46.6
        & One shared key--value head \\
        GQA-8-XXL
        & 0.28
        & 47.1
        & Eight shared key--value groups \\
        \bottomrule
    \end{tabularx}
    \normalsize
\end{table}

GQA is therefore a compact-attention strategy: it reduces the amount of
key--value information stored in the cache through sharing, while retaining
the dense query-to-key interaction pattern. It should not be confused
with sparse attention, which omits selected query--key interactions, or
with linear attention, which changes the mathematical formulation of
attention \parencite{ainslie2023gqa,sun2026efficientattention}.

\subsection{FlashAttention and FlashAttention-2}
\label{subsec:attn-flashattention}

The preceding methods address the cost of retaining key--value
representations during inference. A different problem occurs when the
attention score and probability matrices are materialized during training
or prompt processing. FlashAttention addresses this problem without
approximating the dense attention result. Its central idea is to reorganize
the computation to reduce data movement between high-bandwidth memory
(HBM) and fast on-chip storage.

\paragraph{Memory hierarchy and the I/O bottleneck.}
GPUs contain several levels of memory with different capacities and
bandwidths. HBM provides comparatively large storage, while on-chip SRAM
is smaller and faster. The FlashAttention-2 paper uses NVIDIA A100
hardware as an example and reports HBM bandwidth of approximately
\(1.5\)--\(2.0\) TB/s and an estimated on-chip SRAM bandwidth of
approximately \(19\) TB/s. It also describes approximately \(192\) KB
of SRAM per streaming multiprocessor in the hardware setting considered
by the paper. These figures are hardware-specific examples, but they
illustrate why repeated transfers of large intermediate matrices can
limit practical attention speed \parencite{dao2023flashattention2}.

A conventional implementation may compute \(S=QK^\top\), write \(S\)
to HBM, read it for softmax, write \(P\) to HBM, and read \(P\) again
for \(O=PV\). The score and probability matrices each have
\(N^2\) entries per head. The repeated writes and reads create substantial
memory traffic, and the probability matrix may also need to be retained
for backpropagation. The runtime is consequently influenced by both
arithmetic and memory movement.

\paragraph{Tiling.}
FlashAttention uses tiling to keep intermediate blocks in on-chip SRAM.
The query matrix is divided into row blocks, and the key and value
matrices are divided into corresponding blocks. For a query block
\(Q_i\) and key block \(K_j\), the local score block is
\begin{equation}
    S_{ij}
    =
    \frac{Q_iK_j^\top}{\sqrt{d_k}},
    \qquad
    S_{ij}\in\mathbb{R}^{B_r\times B_c},
    \label{eq:attn-flash-score-tile}
\end{equation}
where \(B_r\) and \(B_c\) are the query-row and key-column block sizes.
The corresponding value block is
\(V_j\in\mathbb{R}^{B_c\times d_v}\). The algorithm loads blocks from
HBM, computes their contribution in SRAM, updates a running output, and
moves to the next block. It avoids writing the full \(N\times N\) score
and probability matrices to HBM.

The blocks cannot simply be softmax-normalized independently and then
concatenated. The softmax denominator for a query row depends on every
permitted key position, including keys in blocks that have not yet been
processed. FlashAttention resolves this dependency through online
softmax.

\paragraph{Online softmax derivation.}
Consider one query-row block and suppose that key--value blocks are
processed in sequence. For a given query row, maintain a running maximum
\(m\), a running normalization denominator \(\ell\), and an
unnormalized value accumulator \(U\). Before processing any key block,
initialize
\begin{equation}
    m^{(0)}=-\infty,
    \qquad
    \ell^{(0)}=0,
    \qquad
    U^{(0)}=0.
    \label{eq:attn-online-softmax-initialization}
\end{equation}
For the next score block \(S^{(j)}\), update the running maximum:
\begin{equation}
    m^{(j)}
    =
    \max\left(
        m^{(j-1)},
        \operatorname{rowmax}(S^{(j)})
    \right).
    \label{eq:attn-online-softmax-max}
\end{equation}
The denominator is updated by rescaling the previous contribution to the
new maximum and adding the current block's contribution:
\begin{equation}
    \boxed{
    \ell^{(j)}
    =
    e^{m^{(j-1)}-m^{(j)}}\ell^{(j-1)}
    +
    \operatorname{rowsum}
    \left(
        e^{S^{(j)}-m^{(j)}}
    \right)
    }.
    \label{eq:attn-online-softmax-denominator}
\end{equation}
The exponentials, maxima, and row sums are evaluated row-wise. The
unnormalized value accumulator is updated as
\begin{equation}
    \boxed{
    U^{(j)}
    =
    e^{m^{(j-1)}-m^{(j)}}U^{(j-1)}
    +
    e^{S^{(j)}-m^{(j)}}V_j
    }.
    \label{eq:attn-online-softmax-output}
\end{equation}
The rescaling factor is necessary because the maximum used to stabilize
the exponential may change when a new block is introduced. After all
key--value blocks have been processed, the exact normalized output is
\begin{equation}
    O=\frac{U^{(T_c)}}{\ell^{(T_c)}},
    \label{eq:attn-online-softmax-final}
\end{equation}
where \(T_c\) is the number of key--value blocks. The division is
performed row-wise and broadcast across the value dimension.

These updates are algebraically equivalent to full softmax attention.
They change the order in which intermediate quantities are computed and
stored, not the attention function itself. The running maximum also
preserves the numerical-stability benefit of max-shifted softmax.
\cref{fig:attn-flash-tiling} illustrates the blockwise computation and
the compact online-softmax state
\parencite{dao2023flashattention2}.

\begin{figure}[htbp]
    \centering
    % Figure source fragment: tiled FlashAttention with online softmax.
% This file is inserted inside a figure environment using \input.
\begin{center}
\setlength{\fboxsep}{5pt}
\renewcommand{\arraystretch}{1.2}
\begin{tabular}{cc}
    \fbox{\parbox[c][0.9cm][c]{0.32\textwidth}{\centering
        \textbf{Query and key tiles}\\
        \(Q_i,\;K_j\)}} &
    \fbox{\parbox[c][0.9cm][c]{0.32\textwidth}{\centering
        \textbf{Value tile}\\
        \(V_j\)}}
\end{tabular}\\[0.3em]
\(\downarrow\hspace{2.8cm}\downarrow\)\\[0.2em]
\begin{tabular}{cc}
    \fbox{\parbox[c][0.9cm][c]{0.32\textwidth}{\centering
        \textbf{Local score tile}\\
        \(S_{ij}=Q_iK_j^\top/\sqrt{d_k}\)}} &
    \fbox{\parbox[c][0.9cm][c]{0.32\textwidth}{\centering
        \textbf{Value tile}\\
        \(V_j\)}}
\end{tabular}\\[0.3em]
\(\searrow\hspace{0.8cm}\swarrow\)\\[0.2em]
\fbox{\parbox[c][1.95cm][c]{0.72\textwidth}{\centering
    \textbf{Online softmax state}\\[0.2em]
    Running maximum \(m\), normalization denominator \(\ell\),\\
    and value accumulator \(U\).\\
    Update the state for each successive key--value tile.}}
\\[0.3em]
\(\downarrow\)\\[0.2em]
\fbox{\parbox[c][0.85cm][c]{0.62\textwidth}{\centering
    \textbf{Normalized output tile}\\
    \(O_i=U/\ell\)}}
\end{center}

    \caption{Conceptual blockwise attention computation. Query, key, and
    value tiles are loaded into on-chip memory; each score tile updates
    the running maximum, normalization denominator, and value accumulator.
    The full score and probability matrices are not materialized in HBM.
    Based on \textcite{dao2023flashattention2}.}
    \label{fig:attn-flash-tiling}
\end{figure}
% Alt text for accessibility register: Query, key, and value matrices are
% shown as blocks. A query block interacts with successive key and value
% blocks. Each interaction updates running maximum m, denominator l, and
% output accumulator U; the final output is U divided by l.

\paragraph{Backward pass and recomputation.}
A conventional backward pass uses the attention probabilities to calculate
gradients. FlashAttention avoids storing the complete probability matrix
in HBM. Instead, it retains compact normalization information and
recomputes the necessary score and probability blocks during the backward
pass. This trades additional arithmetic for reduced memory traffic and
intermediate storage. The gradients are still those of the exact attention
operation:
\begin{align}
    dV&=P^\top dO,
    \label{eq:attn-flash-backward-dv}\\
    dP&=dOV^\top,
    \label{eq:attn-flash-backward-dp}\\
    dS&=\operatorname{SoftmaxBackward}(P,dP),
    \label{eq:attn-flash-backward-ds}\\
    dQ&=\frac{dSK}{\sqrt{d_k}},
    \qquad
    dK=\frac{dS^\top Q}{\sqrt{d_k}}.
    \label{eq:attn-flash-backward-dqdk}
\end{align}
The algorithm recomputes local blocks when required rather than retaining
all \(N\times N\) intermediates. The additional memory required by the
attention operation is thereby reduced from quadratic to linear in
sequence length, apart from the storage of the input and output tensors
\parencite{dao2023flashattention2}.

\paragraph{FlashAttention-2 and GPU work partitioning.}
FlashAttention-2 improves the work distribution of the original
algorithm. The paper identifies three related changes. It reduces
non-matrix-multiplication operations, parallelizes attention across
sequence blocks even when a single head is being processed, and changes
the division of work among warps within a thread block to reduce
communication through shared memory.

The distinction between matrix and non-matrix operations matters because
modern GPUs provide specialized matrix-multiplication hardware. The
FlashAttention-2 paper gives an A100 example in which the reported
theoretical throughput for FP16/BF16 matrix multiplication is
\(312\) TFLOPs/s, compared with \(19.5\) TFLOPs/s for non-matrix
FP32 operations. The ratio is
\begin{equation}
    \frac{312}{19.5}=16.
    \label{eq:attn-a100-throughput-ratio}
\end{equation}
This comparison is specific to the hardware and precisions reported in
the paper. It illustrates why a kernel that reduces non-matrix work and
increases the fraction of time spent in optimized matrix operations can
achieve better hardware utilization.

The original FlashAttention work distributes computation across batch
and head dimensions. FlashAttention-2 additionally partitions work along
the sequence dimension, increasing the number of independent thread
blocks when batch size or head count is small. Within a thread block, it
also changes the warp-level partitioning to reduce shared-memory
communication and synchronization. For causal attention, blocks entirely
outside the permitted triangular region can be skipped rather than
computed and discarded.

The paper reports approximately a twofold speedup over the original
FlashAttention in its evaluated settings, with \(50\)--\(73\%\) of
the theoretical maximum FLOPs/s on A100. In end-to-end GPT-style
training, it reports up to \(225\) TFLOPs/s per A100 GPU and
approximately \(72\%\) model FLOPs utilization. These are experimental
results tied to the implementations and hardware in the paper; they
should not be interpreted as fixed speedups for every model or device
\parencite{dao2023flashattention2}.

The central distinction is that FlashAttention-2 does not remove the
quadratic arithmetic complexity of dense attention. It changes the
execution schedule, intermediate storage, and work partitioning so that
the same mathematical operation is performed with less data movement
and better hardware utilization.

\subsection{Linear attention}
\label{subsec:attn-linear-attention}

FlashAttention improves the execution of dense attention without changing
its mathematical result. Linear attention takes a different approach:
it changes the formulation so that the sequence-length dependence can be
linear rather than quadratic. The supplied survey by
\textcite{sun2026efficientattention} organizes linear attention into
kernelized methods, recurrent methods with forgetting mechanisms, and
methods that interpret the recurrent state as a fast-weight or
in-context-learning mechanism.

\paragraph{Kernelized reformulation.}
Consider one query \(q_i\) and the sequence of keys and values.
In the efficient-attention survey's notation, the softmax kernel
contributes weights proportional to \(\exp(q_i k_j^\top)\). Kernelized linear attention seeks a feature
mapping
\begin{equation}
    \phi:\mathbb{R}^{d_k}\longrightarrow\mathbb{R}^{r}
    \label{eq:attn-linear-feature-map}
\end{equation}
such that
\begin{equation}
    \exp(q_i k_j^\top)
    \approx
    \phi(q_i)^\top\phi(k_j).
    \label{eq:attn-linear-kernel-approximation}
\end{equation}
If scaled dot products are used, the scaling can be incorporated into
the arguments of the feature map. The essential requirement is that
the similarity kernel be approximated by an inner product in a feature
space.

Substituting the feature approximation into normalized attention gives
\begin{equation}
    o_i
    \approx
    \frac{
        \sum_{j=1}^{N}
        \phi(q_i)^\top\phi(k_j)v_j
    }{
        \sum_{j=1}^{N}
        \phi(q_i)^\top\phi(k_j)
    }.
    \label{eq:attn-linear-substitution}
\end{equation}
Since \(\phi(q_i)\) does not depend on \(j\), associativity and
distributivity allow it to be factored out:
\begin{align}
    o_i
    &\approx
    \frac{
        \phi(q_i)^\top
        \left(
            \sum_{j=1}^{N}\phi(k_j)v_j^\top
        \right)
    }{
        \phi(q_i)^\top
        \left(
            \sum_{j=1}^{N}\phi(k_j)
        \right)
    }.
    \label{eq:attn-linear-reordered}
\end{align}
Define the aggregated key--value matrix and key sum as
\begin{align}
    A&=
    \sum_{j=1}^{N}\phi(k_j)v_j^\top,
    & A&\in\mathbb{R}^{r\times d_v},
    \label{eq:attn-linear-state-matrix}\\
    z&=
    \sum_{j=1}^{N}\phi(k_j),
    & z&\in\mathbb{R}^{r}.
    \label{eq:attn-linear-normalizer-state}
\end{align}
Then
\begin{equation}
    \boxed{
    o_i
    \approx
    \frac{\phi(q_i)^\top A}
         {\phi(q_i)^\top z}
    }.
    \label{eq:attn-linear-output-state}
\end{equation}
The numerator and denominator are evaluated using the aggregated state.
The full \(N\times N\) attention matrix is not constructed.

For the complete sequence, let
\begin{equation}
    \Phi_Q=\phi(Q)\in\mathbb{R}^{N\times r},
    \qquad
    \Phi_K=\phi(K)\in\mathbb{R}^{N\times r}.
    \label{eq:attn-linear-feature-matrices}
\end{equation}
The output can be written as
\begin{equation}
    \boxed{
    O
    =
    \frac{
        \Phi_Q(\Phi_K^\top V)
    }{
        \Phi_Q(\Phi_K^\top\mathbf{1}_N)
    }
    }.
    \label{eq:attn-linear-matrix-form}
\end{equation}
Here, \(\mathbf{1}_N\) is an \(N\)-dimensional vector of ones, and the
division is row-wise, with the denominator broadcast over the value
dimension. The matrix product \(\Phi_K^\top V\) has shape
\(r\times d_v\), while \(\Phi_Q(\Phi_K^\top V)\) has shape
\(N\times d_v\).

The first product costs
\(\mathcal{O}(Nrd_v)\), and the multiplication by \(\Phi_Q\) costs
another \(\mathcal{O}(Nrd_v)\). The denominator costs
\(\mathcal{O}(Nr)\). The total is therefore
\begin{equation}
    \mathcal{O}(Nrd_v+Nr),
    \label{eq:attn-linear-complexity}
\end{equation}
which is linear in \(N\) when the feature and value dimensions are fixed.
If \(r=d_v=d_h\), the leading cost is
\(\mathcal{O}(Nd_h^2)\). The reduction in sequence-length complexity
comes from changing the order of matrix multiplication and avoiding
the \(N\times N\) interaction matrix
\parencite{sun2026efficientattention}.

\paragraph{Feature-map choices and approximation.}
The feature map determines how well the linearized computation
approximates the original softmax kernel. The survey describes Linear
Transformer as using the positive feature map
\begin{equation}
    \phi(x)=\operatorname{ELU}(x)+1,
    \label{eq:attn-linear-elu-feature}
\end{equation}
where \(\operatorname{ELU}\) is the exponential linear unit. The positive
outputs are compatible with the non-negative character of the
exponential kernel and permit a normalization term.

Performer uses FAVOR+, a positive orthogonal random-feature method
designed to provide an unbiased approximation of the softmax kernel.
The survey also discusses Random Feature Attention and cosFormer as
alternative feature-based constructions. These methods differ in how
they represent or approximate the similarity function. Their computational
advantage is obtained by replacing exact softmax attention with a
feature-space approximation, so approximation quality and the feature
dimension are part of the design problem
\parencite{sun2026efficientattention}.

\paragraph{Recurrent state formulation.}
The associative form can be evaluated incrementally. Define the state
after token \(t\) as
\begin{equation}
    A_t
    =
    \sum_{j=1}^{t}\phi(k_j)v_j^\top,
    \qquad
    z_t
    =
    \sum_{j=1}^{t}\phi(k_j).
    \label{eq:attn-linear-recurrent-state}
\end{equation}
The state updates are
\begin{align}
    A_t&=A_{t-1}+\phi(k_t)v_t^\top,
    \label{eq:attn-linear-state-update}\\
    z_t&=z_{t-1}+\phi(k_t),
    \label{eq:attn-linear-normalizer-update}
\end{align}
with \(A_0=0\) and \(z_0=0\). The current output is
\begin{equation}
    o_t
    =
    \frac{\phi(q_t)^\top A_t}
         {\phi(q_t)^\top z_t}.
    \label{eq:attn-linear-recurrent-output}
\end{equation}
The accumulated state summarizes the preceding key--value information.
The current query retrieves from that state without separately comparing
it with every previously stored key.

The state has fixed dimensions \(r\times d_v\) and \(r\), so the update
and storage per token do not grow with the number of preceding tokens,
provided \(r\) and \(d_v\) are fixed. This gives linear-time sequence
processing and constant-size recurrent memory per step. The recurrent
form is useful for autoregressive decoding, while parallel or chunkwise
forms can be used to expose more parallelism during training and prompt
processing \parencite{sun2026efficientattention}.

\paragraph{Forgetting mechanisms.}
Basic additive linear attention accumulates previous information without
an explicit decay. The survey describes recurrent designs that introduce
forgetting mechanisms so that the contribution of older information can
decrease over time.

RetNet is presented as a data-independent-decay example. Its state
recurrence is
\begin{equation}
    S_t=\gamma S_{t-1}+k_t^\top v_t,
    \qquad
    0<\gamma<1,
    \label{eq:attn-retnet-recurrence}
\end{equation}
where \(\gamma\) is a learned decay coefficient, commonly associated
with a retention head. The output is
\begin{equation}
    o_t=q_tS_t.
    \label{eq:attn-retnet-output}
\end{equation}
Unrolling the recurrence gives
\begin{equation}
    S_t
    =
    \sum_{n=1}^{t}
    \gamma^{t-n}k_n^\top v_n,
    \qquad
    o_t
    =
    \sum_{n=1}^{t}
    \gamma^{t-n}q_tk_n^\top v_n.
    \label{eq:attn-retnet-unrolled}
\end{equation}
The contribution from token \(n\) is multiplied by
\(\gamma^{t-n}\). When \(\gamma\) is closer to one, information decays
more slowly; a smaller \(\gamma\) produces faster forgetting. The decay
is data-independent because the coefficient is not recalculated from
the content of each incoming token.

Data-dependent approaches instead use a gate \(G_t\) determined by the
current input:
\begin{equation}
    S_t=G_tS_{t-1}+k_t^\top v_t.
    \label{eq:attn-data-dependent-gate}
\end{equation}
A gate value near zero suppresses the corresponding part of the previous
state, while a value near one preserves it. The survey discusses Mamba,
Gated Linear Attention, xLSTM, GateLoop, HGRN, and related architectures
as examples of input-dependent or gated recurrent state dynamics.
The exact structure of \(G_t\) varies across architectures; it should
not be assumed to be a single scalar in every model
\parencite{sun2026efficientattention}.

\paragraph{Linear attention as an in-context memory.}
The recurrent state can also be interpreted as an associative memory.
The update \(A_t=A_{t-1}+k_t^\top v_t\) writes a key--value association
into a matrix-valued state, and \(o_t=q_tA_t\) retrieves information
from that state using the current query. The state is therefore both
a compact summary of preceding tokens and a rapidly updated memory.
The efficient-attention survey places fast-weight and test-time-regression
approaches, including DeltaNet and related methods, within this broader
line of work. Their state-update rules modify the simple additive
association update to control how new information changes previously
stored information \parencite{sun2026efficientattention}.

Linear attention changes the computation rather than merely changing
the implementation schedule. Its sequence-length complexity can be
linear, but its output is generally an approximation or alternative
formulation of softmax attention. The feature map, state dynamics,
forgetting mechanism, and ability to retain detailed long-range
information therefore affect its modeling behavior.

\subsection{Sparse attention}
\label{subsec:attn-sparse-attention}

Linear attention changes the kernel or recurrent formulation. Sparse
attention instead retains softmax-style interactions but restricts the
set of key positions available to each query. Its motivation is that
many query--key interactions may be unnecessary for a useful output.
The survey by \textcite{sun2026efficientattention} organizes sparse
attention into fixed-pattern attention, block-sparse attention,
clustering-based attention, and related bidirectional designs.

\paragraph{Sparse interaction sets.}
Let \(\mathcal{S}(i)\) be the subset of key positions selected for query
\(i\). Sparse attention is
\begin{equation}
    o_i
    =
    \sum_{j\in\mathcal{S}(i)}
    \alpha_{ij}v_j,
    \label{eq:attn-sparse-output}
\end{equation}
where
\begin{equation}
    \alpha_{ij}
    =
    \frac{\exp(s_{ij})}
         {\displaystyle\sum_{r\in\mathcal{S}(i)}\exp(s_{ir})},
    \qquad
    s_{ij}=\frac{q_i k_j^\top}{\sqrt{d_k}}.
    \label{eq:attn-sparse-softmax}
\end{equation}
The softmax denominator includes only the selected positions. If each
query attends to at most \(k\) keys, the main interaction cost becomes
\begin{equation}
    \mathcal{O}(Nkd_h)
    \label{eq:attn-sparse-complexity}
\end{equation}
for one head. When \(k\) is fixed as \(N\) grows, this cost is linear
in sequence length. The quality of the result depends on whether the
selected set preserves the information needed by the query.

\paragraph{Fixed-pattern sparse attention.}
A fixed-pattern method defines the permitted interactions before
evaluating the input. Local-window attention is one example. In a
causal window of width \(w\), query position \(i\) attends to recent
positions:
\begin{equation}
    \mathcal{S}_{\mathrm{local}}(i)
    =
    \left\{
        j:
        \max(1,i-w+1)\leq j\leq i
    \right\}.
    \label{eq:attn-local-window}
\end{equation}
The interaction count is at most \(Nw\), so the leading cost is
\(\mathcal{O}(Nwd_h)\). A fixed window is efficient when \(w\) is
small, but a query cannot directly retrieve a token outside the window.

Dilated attention selects positions separated by a fixed interval.
For dilation rate \(r\), an illustrative causal set is
\begin{equation}
    \mathcal{S}_{\mathrm{dilated}}(i)
    =
    \{i,i-r,i-2r,\ldots\}\cap\{1,\ldots,N\}.
    \label{eq:attn-dilated-pattern}
\end{equation}
This pattern expands the temporal or positional range without attending
to every intervening token. The survey lists local-window and dilated
designs among fixed sparse attention methods
\parencite{sun2026efficientattention}.

\paragraph{Block-sparse attention.}
Token-level sparsity can create irregular memory accesses. Block-sparse
attention groups tokens into blocks and selects interactions between
blocks. Let the block size be \(b\), and let
\(T=\lceil N/b\rceil\) be the number of blocks. A block mask
\begin{equation}
    M_{\mathrm{block}}\in\{0,1\}^{T\times T}
    \label{eq:attn-block-mask}
\end{equation}
indicates which query blocks may interact with which key blocks. If each
query block attends to at most \(k_b\) key blocks, then the number of
token-level interactions is proportional to
\[
    \frac{N}{b}\,k_b b^2
    =
    Nk_b b.
\]
The leading attention cost is therefore
\begin{equation}
    \mathcal{O}(Nk_bbd_h).
    \label{eq:attn-block-sparse-complexity}
\end{equation}
For fixed \(b\) and \(k_b\), the dependence on \(N\) is linear.

Block structure allows hardware kernels to skip entire
matrix-multiplication tiles. The efficient-attention survey distinguishes
block-sparse methods for prompt prefill from methods designed for decoding. Prefill methods include approaches
that select important blocks while processing a long input; decoding
methods select blocks from the existing KV cache for each new query.
The survey also discusses routing-based block sparsity, including
methods that use pooled block representations or learned selection
scores \parencite{sun2026efficientattention}.

\paragraph{Clustering and retrieval-based attention.}
Clustering-based methods organize key representations into groups and
retrieve selected groups rather than comparing every query with every
key. Suppose the key vectors are partitioned into \(C\) clusters with
centroids
\begin{equation}
    \mu_c
    =
    \frac{1}{|\mathcal{C}_c|}
    \sum_{j\in\mathcal{C}_c}k_j,
    \qquad c=1,\ldots,C.
    \label{eq:attn-cluster-centroid}
\end{equation}
A query can be compared with the centroids to identify promising
clusters, after which the corresponding token-level keys and values
are retrieved. The survey describes ClusterKV as using \(k\)-means
clustering and RetrievalAttention as using approximate nearest-neighbor
search. It also discusses Locality-Sensitive Hashing (LSH) as a way
to map similar query and key vectors into common hash buckets
\parencite{sun2026efficientattention}.

The efficiency of retrieval-based attention depends on the quality of
the index or clustering structure and on the cost of selecting candidate
groups. If relevant keys are assigned to clusters that are not retrieved,
the query can lose access to useful information. The retrieval stage
must therefore be included when evaluating the complete computational
cost.

\paragraph{Sparsity and hardware alignment.}
A sparse mathematical pattern does not automatically guarantee a
proportional wall-clock speedup. Irregular accesses, selection overhead,
small matrix tiles, and insufficient parallel work can reduce the benefit
of sparsity. Block-sparse designs attempt to align the pattern with
matrix-multiplication tiles so that skipped interactions correspond to
operations the hardware can omit efficiently. The survey emphasizes
that practical sparse attention must account for block sizes, grouping
of key--value heads, and reuse of selected blocks
\parencite{sun2026efficientattention}.

The survey also identifies training sparse attention from scratch as
a challenge. Restricted connectivity can lead to sparse gradient
signals and less effective optimization. Training-aware or trainable
sparse mechanisms attempt to learn useful selection patterns, but
high sparsity can still harm model quality. Sparse attention is therefore
a trade-off between omitted computation and the risk of excluding
information needed for accurate prediction
\parencite{sun2026efficientattention}.

\subsection{Hybrid attention architectures}
\label{subsec:attn-hybrid-architectures}

Dense, compact, hardware-efficient, sparse, and linear attention
address different constraints. Dense attention offers unrestricted
token-to-token interaction. Compact methods reduce the amount of
key--value state that must be retained. Hardware-efficient methods
reduce data movement and improve the execution of the same mathematical
operation. Sparse methods omit selected interactions, while linear
methods reformulate the computation to avoid the full pairwise matrix.
A hybrid architecture combines more than one of these mechanisms so
that different layers or components can provide complementary
capabilities.

The efficient-attention survey describes several broad hybrid patterns.
A sparse--dense design interleaves dense attention layers with locally
banded or otherwise sparse layers. Dense layers provide global
interactions, while sparse layers reduce the number of token pairs
processed in other parts of the network. The survey discusses GPT-3
as an example of interleaving dense and locally banded sparse attention
layers \parencite{sun2026efficientattention}.

A linear--full design uses linear or recurrent attention in many layers
and inserts softmax-based full attention at selected depths. The survey
discusses Jamba and MiniMax-01 as examples of this pattern. In the
reported MiniMax-01 design, Lightning Attention is used across most
layers, with full attention inserted periodically. The purpose is to
retain efficient state-based processing while preserving opportunities
for precise global softmax interactions
\parencite{sun2026efficientattention}.

A local--full design alternates sliding-window attention with occasional
global attention layers. The survey discusses Gemma 3, Command A, and
LLaMA-4 Maverick as examples of architectures using local and global
attention components. Local layers emphasize nearby context, while
global layers provide longer-range interactions. The frequency and
positioning of global layers determine the balance between computation
and global information exchange \parencite{sun2026efficientattention}.

These examples demonstrate that efficient attention is not necessarily
a choice between one universal mechanism and another. Different
architectural components can be combined to distribute the requirements
of local processing, global retrieval, recurrent memory, and exact
softmax interactions across the network.

\begin{figure}[htbp]
    \centering
    % Figure source fragment: taxonomy of efficient attention mechanisms.
% This file is inserted inside a figure environment using \input.
\begin{center}
\setlength{\fboxsep}{5pt}
\renewcommand{\arraystretch}{1.2}
\begin{tabular}{cc}
    \fbox{\parbox[t][1.8cm][t]{0.43\textwidth}{\centering
        \textbf{Hardware-efficient attention}\\[0.2em]
        \scriptsize Same attention function; reduce data movement and
        improve kernel execution.\\
        Example: FlashAttention-2.}} &
    \fbox{\parbox[t][1.8cm][t]{0.43\textwidth}{\centering
        \textbf{Compact attention}\\[0.2em]
        \scriptsize Reduce stored key--value state through sharing or
        compression.\\
        Examples: MQA and GQA.}}\\[0.35em]
    \fbox{\parbox[t][1.8cm][t]{0.43\textwidth}{\centering
        \textbf{Sparse attention}\\[0.2em]
        \scriptsize Select a subset of token or block interactions.\\
        Examples: local, block-sparse, and retrieval-based attention.}} &
    \fbox{\parbox[t][1.8cm][t]{0.43\textwidth}{\centering
        \textbf{Linear attention}\\[0.2em]
        \scriptsize Reformulate similarity computation or use recurrent
        state updates.\\
        Examples: kernelized and gated linear attention.}}
\end{tabular}

\par\smallskip
\footnotesize
Hybrid architectures combine multiple mechanisms across layers or components.
\normalsize
\end{center}

    \caption{Conceptual taxonomy of efficient attention mechanisms.
    Hardware-efficient attention changes the execution schedule and
    memory traffic; compact attention reduces stored key--value state;
    sparse attention selects a subset of interactions; and linear
    attention changes the mathematical formulation to obtain
    subquadratic sequence-length dependence. Hybrid architectures
    combine these mechanisms. Based on
    \textcite{sun2026efficientattention}.}
    \label{fig:attn-efficiency-taxonomy}
\end{figure}
% Alt text for accessibility register: Four categories are displayed:
% hardware-efficient attention, compact attention, sparse attention, and
% linear attention. Each category is paired with its principal mechanism.
% A final note indicates that hybrid architectures combine categories.

\subsection{Limitations and connection to the next topic}
\label{subsec:attn-limitations-transition}

Attention provides content-dependent information exchange in the
Transformer. Its mathematical structure separates compatibility
estimation from value aggregation, while multi-head attention provides
several learned interaction spaces. Causal masking constrains information
flow for autoregressive prediction, and backpropagation optimizes the
projection spaces through the loss function.

The architectural and computational variants address different
limitations and should not be treated as interchangeable. Multi-query
and grouped-query attention reduce key--value cache storage and memory
traffic, but retain dense query-to-key interactions. FlashAttention-2
preserves exact dense attention and its quadratic arithmetic dependence
while reducing intermediate memory traffic and improving GPU work
partitioning. Kernelized linear attention changes the similarity
formulation and introduces approximation or alternative recurrent
state dynamics. Sparse attention omits selected interactions and
depends on the quality and hardware efficiency of its selection
pattern. Hybrid architectures combine these mechanisms, but their
benefits depend on the layer arrangement, sequence length, task, and
hardware environment.

The central engineering lesson is that arithmetic complexity, memory
capacity, memory bandwidth, approximation error, and hardware
parallelism are distinct quantities. A reduction in one does not
guarantee a reduction in all the others. Consequently, an attention
architecture must be evaluated in relation to its mathematical
formulation and its intended training or inference workload.

This completes the architectural account of how Transformer layers
exchange contextual information. The next chapter turns from the
architecture of the computation to the optimization procedures used
to learn its parameters. Understanding attention's forward and
backward computations provides the necessary basis for examining
optimization algorithms, loss functions, and the numerical and
computational techniques used to train large language models.
